
\documentclass{article}
\usepackage{colortbl}
\usepackage{makecell}
\usepackage{multirow}
\usepackage{supertabular}

\begin{document}

\newcounter{utterance}

\centering \large Interaction Transcript for game `hot\_air\_balloon', experiment `air\_balloon\_survival\_mt\_negotiation\_easy', episode 2 with mistral{-}large{-}3{-}azure.
\vspace{24pt}

{ \footnotesize  \setcounter{utterance}{1}
\setlength{\tabcolsep}{0pt}
\begin{supertabular}{c@{$\;$}|p{.15\linewidth}@{}p{.15\linewidth}p{.15\linewidth}p{.15\linewidth}p{.15\linewidth}p{.15\linewidth}}
    \# & $\;$A & \multicolumn{4}{c}{Game Master} & $\;\:$B\\
    \hline

    \theutterance \stepcounter{utterance}  
    & & \multicolumn{4}{p{0.6\linewidth}}{
        \cellcolor[rgb]{0.9,0.9,0.9}{
            \makecell[{{p{\linewidth}}}]{
                \texttt{\tiny{[P1$\langle$GM]}}
                \texttt{Int qiegħed tipparteċipa f'logħba ta' negozjar kollaborattiva.} \\
\\ 
\texttt{Flimkien ma' parteċipant ieħor, trid taqblu fuq sett wieħed ta' oġġetti li se jinżammu. Kull wieħed minnkom għandu l{-}fehma tiegħu dwar kemm kull oġġett hu importanti għalih (importanza). Int ma tafx kif il{-}parteċipant l{-}ieħor jivvaluta l{-}oġġetti. Barra minn hekk, tingħata l{-}isforz li jitlob kull oġġett.} \\
\texttt{Tistgħu taqblu fuq sett biss jekk l{-}isforz totali tal{-}oġġetti magħżula ma jaqbiżx limitu kondiviż:} \\
\\ 
\texttt{LIMITU = 3288} \\
\\ 
\texttt{Hawn huma l{-}valuri individwali tal{-}isforz tal{-}oġġetti:} \\
\\ 
\texttt{Isforz tal{-}oġġett = \{"C08": 329, "C55": 41, "A43": 417, "A67": 75, "B42": 390, "A62": 887, "B19": 808, "C43": 151, "C73": 849, "C10": 129, "B36": 350, "A18": 118, "C61": 630, "C67": 602, "C42": 801\}} \\
\\ 
\texttt{Hawn hi l{-}fehma personali tiegħek dwar l{-}importanza ta' kull oġġett:} \\
\\ 
\texttt{Valuri tal{-}importanza tal{-}oġġetti = \{"C08": 138, "C55": 583, "A43": 868, "A67": 822, "B42": 783, "A62": 65, "B19": 262, "C43": 121, "C73": 508, "C10": 780, "B36": 461, "A18": 484, "C61": 668, "C67": 389, "C42": 808\}} \\
\\ 
\texttt{Għan:} \\
\\ 
\texttt{L{-}għan tiegħek hu li tinnegozja sett kondiviż ta' oġġetti li jibbenefikak kemm jista' jkun (jiġifieri, jimmassimizza l{-}importanza totali GĦALIK), filwaqt li tibqa' fil{-}LIMITU. M'intix obbligat tagħmel PROPOSTA f'kull messaġġ {-} tista' sempliċement tinnegozja wkoll. It{-}tattiċi kollha huma permessi!} \\
\\ 
\texttt{Protokoll ta' Interazzjoni:} \\
\\ 
\texttt{Tista' tuża biss il{-}formati strutturati li ġejjin f'messaġġ:} \\
\\ 
\texttt{PROPOSTA: \{'A', 'B', 'C', …\}} \\
\texttt{Ipproponi li jinżammu eżattament dawk l{-}oġġetti.} \\
\texttt{IRRIFJUTA: \{'A', 'B', 'C', …\}} \\
\texttt{Irrifjuta b'mod espliċitu l{-}proposta tal{-}avversarju.} \\
\texttt{ARGUMENT: \{'...'\}} \\
\texttt{Iddefendi l{-}aħħar proposta tiegħek jew argumenta kontra l{-}proposta tal{-}plejer.} \\
\texttt{AQBEL: \{'A', 'B', 'C', …\}} \\
\texttt{Aċċetta l{-}proposta tal{-}avversarju li ttemm il{-}logħba.} \\
\\ 
\\ 
\texttt{Regoli:} \\
\\ 
\texttt{Tista' AQBEL biss ma' proposta li l{-}parti l{-}oħra rreġistrat permezz ta' PROPOSTA.} \\
\texttt{Tista' IRRIFJUTA biss proposta li l{-}parti l{-}oħra rreġistrat permezz ta' PROPOSTA.} \\
\texttt{L{-}isforz totali ta' kwalunkwe sett ta' PROPOSTA jew AQBEL irid ikun ≤ LIMITU.} \\
\texttt{TIŻVELAX il{-}punteġġi moħbija tal{-}importanza tiegħek.} \\
\texttt{Tikketta f'format strutturat trid tkun segwita minn kolon u spazju vojt. L{-}argument irid ikun sett ta' python li fih 0 jew aktar strings.} \\
\texttt{Għalhekk, irid ikun fil{-}forma "TIKKETTA: \{...\}"} \\
\texttt{Segwi b'mod strett il{-}protokoll ta' interazzjoni u TIKTIBX xejn lil hinn mill{-}istruttura mogħtija.} \\
\texttt{Il{-}logħba tispiċċa meta naħa waħda tagħti AQBEL għal PROPOSTA magħmula mill{-}plejer l{-}ieħor.} \\
\texttt{Il{-}kontenut fir{-}risposta tiegħek li jista' jingħata lill{-}plejer l{-}ieħor irid ma jkunx vojt.} \\
\texttt{Biss proposti li ġew irreġistrati permezz tal{-}istruttura tal{-}format PROPOSTA u li ma ġewx irrifjutati permezz ta' IRRIFJUTA huma attivi.} \\
\\ 
            }
        }
    }
    & \\ \\

    \theutterance \stepcounter{utterance}  
    & \multicolumn{4}{p{0.6\linewidth}}{
        \cellcolor[rgb]{0.9,0.9,0.9}{
            \makecell[{{p{\linewidth}}}]{
                \texttt{\tiny{[P1$\rangle$GM]}}
                \texttt{PROPOSTA: \{'C55', 'A43', 'A67', 'C43', 'C10', 'A18', 'C67'\}} \\
\\ 
\texttt{(Isforz totali: 41 + 417 + 75 + 151 + 129 + 118 + 602 = 1533 ≤ LIMITU)} \\
\texttt{(Importanza tiegħi: 583 + 868 + 822 + 121 + 780 + 484 + 389 = 4047)} \\
            }
        }
    }
    & & \\ \\

    \theutterance \stepcounter{utterance}  
    & & \multicolumn{4}{p{0.6\linewidth}}{
        \cellcolor[rgb]{0.9,0.9,0.9}{
            \makecell[{{p{\linewidth}}}]{
                \texttt{\tiny{[P1$\langle$GM]}}
                \texttt{Int użajt tikketta PROPOSTA, iżda ma pprovdejtx sett python validu li fih strings bħala argumenti, eż. \{'A', 'B', 'C', ...\}. Erġa' pprova.} \\
            }
        }
    }
    & \\ \\

    \theutterance \stepcounter{utterance}  
    & \multicolumn{4}{p{0.6\linewidth}}{
        \cellcolor[rgb]{0.9,0.9,0.9}{
            \makecell[{{p{\linewidth}}}]{
                \texttt{\tiny{[P1$\rangle$GM]}}
                \texttt{PROPOSTA: \{'C55', 'A43', 'A67', 'C43', 'C10', 'A18', 'C67'\}} \\
            }
        }
    }
    & & \\ \\

    \theutterance \stepcounter{utterance}  
    & & \multicolumn{4}{p{0.6\linewidth}}{
        \cellcolor[rgb]{0.9,0.9,0.9}{
            \makecell[{{p{\linewidth}}}]{
                \texttt{\tiny{[GM$\rangle$P2]}}
                \texttt{Qed tipparteċipa f'logħba ta' negozjar kollaborattiva.} \\
\\ 
\texttt{Flimkien ma' parteċipant ieħor, trid taqblu fuq sett wieħed ta' oġġetti li se jinżammu. Kull wieħed minnkom għandu l{-}fehma tiegħu dwar kemm kull oġġett hu importanti għalih (importanza). Ma tafx kif il{-}parteċipant l{-}ieħor jivvaluta l{-}oġġetti. Barra minn hekk, tingħata l{-}isforz li jitlob kull oġġett.} \\
\texttt{Tistgħu taqblu fuq sett biss jekk l{-}isforz totali tal{-}oġġetti magħżula ma jaqbiżx limitu kondiviż:} \\
\\ 
\texttt{LIMITU = 3288} \\
\\ 
\texttt{Hawn huma l{-}valuri individwali tal{-}isforz tal{-}oġġetti:} \\
\\ 
\texttt{Isforz tal{-}oġġett = \{"C08": 329, "C55": 41, "A43": 417, "A67": 75, "B42": 390, "A62": 887, "B19": 808, "C43": 151, "C73": 849, "C10": 129, "B36": 350, "A18": 118, "C61": 630, "C67": 602, "C42": 801\}} \\
\\ 
\texttt{Hawn hi l{-}fehma personali tiegħek dwar l{-}importanza ta' kull oġġett:} \\
\\ 
\texttt{Valuri tal{-}importanza tal{-}oġġetti = \{"C08": 138, "C55": 583, "A43": 868, "A67": 822, "B42": 783, "A62": 65, "B19": 262, "C43": 121, "C73": 508, "C10": 780, "B36": 461, "A18": 484, "C61": 668, "C67": 389, "C42": 808\}} \\
\\ 
\texttt{Għan:} \\
\\ 
\texttt{L{-}għan tiegħek hu li tinnegozja sett kondiviż ta' oġġetti li jibbenefikak kemm jista' jkun (jiġifieri, jimmassimizza l{-}importanza totali GĦALIK), filwaqt li tibqa' fil{-}LIMITU. M'intix obbligat tagħmel PROPOSTA f'kull messaġġ {-} tista' sempliċement tinnegozja wkoll. It{-}tattiċi kollha huma permessi!} \\
\\ 
\texttt{Protokoll ta' Interazzjoni:} \\
\\ 
\texttt{Tista' tuża biss il{-}formati strutturati li ġejjin f'messaġġ:} \\
\\ 
\texttt{PROPOSTA: \{'A', 'B', 'C', …\}} \\
\texttt{Ipproponi li jinżammu eżatt dawk l{-}oġġetti.} \\
\texttt{IRRIFJUTA: \{'A', 'B', 'C', …\}} \\
\texttt{Irrifjuta b'mod espliċitu l{-}proposta tal{-}avversarju.} \\
\texttt{ARGUMENT: \{'...'\}} \\
\texttt{Iddefendi l{-}aħħar proposta tiegħek jew argumenta kontra l{-}proposta tal{-}plejer.} \\
\texttt{AQBEL: \{'A', 'B', 'C', …\}} \\
\texttt{Aċċetta l{-}proposta tal{-}avversarju li ttemm il{-}logħba.} \\
\\ 
\\ 
\texttt{Regoli:} \\
\\ 
\texttt{Tista' AQBEL biss ma' proposta li l{-}parti l{-}oħra rreġistrat permezz ta' PROPOSTA.} \\
\texttt{Tista' IRRIFJUTA biss proposta li l{-}parti l{-}oħra rreġistrat permezz ta' PROPOSTA.} \\
\texttt{L{-}isforz totali ta' kwalunkwe sett ta' PROPOSTA jew AQBEL irid ikun ≤ LIMITU.} \\
\texttt{TIŻVELAX il{-}punteġġi moħbija tal{-}importanza tiegħek.} \\
\texttt{Tikketta f'format strutturat trid tkun segwita b'kolon u spazju. L{-}argument irid ikun sett ta' Python li jkun fih 0 jew aktar strings.} \\
\texttt{Għalhekk, irid ikun fil{-}forma "TIKKETTA: \{...\}"} \\
\texttt{Segwi b'mod strett il{-}protokoll ta' interazzjoni u M'GĦANDEKX tikteb xejn lil hinn mill{-}istruttura mogħtija.} \\
\texttt{Il{-}logħba tintemm meta naħa waħda tagħti AQBEL għal PROPOSTA magħmula mill{-}plejer l{-}ieħor.} \\
\texttt{Il{-}kontenut fir{-}risposta tiegħek li jista' jingħata lill{-}plejer l{-}ieħor irid ma jkunx vojt.} \\
\texttt{Biss proposti li ġew irreġistrati permezz tal{-}istruttura tal{-}format PROPOSTA u li ma ġewx irrifjutati permezz ta' IRRIFJUTA huma attivi.} \\
\\ 
\\ 
\\ 
\texttt{Issa se tirċievi l{-}ewwel messaġġ tal{-}plejer l{-}ieħor.} \\
\\ 
\texttt{PROPOSTA: \{'C55', 'A43', 'A67', 'C43', 'C10', 'A18', 'C67'\}} \\
            }
        }
    }
    & \\ \\

    \theutterance \stepcounter{utterance}  
    & & & \multicolumn{4}{p{0.6\linewidth}}{
        \cellcolor[rgb]{0.9,0.9,0.9}{
            \makecell[{{p{\linewidth}}}]{
                \texttt{\tiny{[GM$\langle$P2]}}
                \texttt{IRRIFJUTA: \{'C55', 'A43', 'A67', 'C43', 'C10', 'A18', 'C67'\}} \\
\\ 
\texttt{ARGUMENT: \{'Din il{-}proposta toħloq isforz totali ta’ 1433, iżda l{-}importanza totali tiegħi għal dawn l{-}oġġetti hija biss 3726, li hija baxxa ħafna. F’dan il{-}kuntest, l{-}oġġetti li joffruni l{-}aħjar ratio ta’ importanza għal isforz huma "B42" (783/390=2.01), "C73" (508/849=0.60), u "B36" (461/350=1.32). Nistgħu nikkonċentraw fuq sett li jinkludi oġġetti b’ratio għoli u jibqa’ fil{-}limitu.'\}} \\
\\ 
\texttt{PROPOSTA: \{'B42', 'C73', 'B36', 'C61', 'C42'\}} \\
            }
        }
    }
     \\ \\

    \theutterance \stepcounter{utterance}  
    & & \multicolumn{4}{p{0.6\linewidth}}{
        \cellcolor[rgb]{0.9,0.9,0.9}{
            \makecell[{{p{\linewidth}}}]{
                \texttt{\tiny{[P1$\langle$GM]}}
                \texttt{IRRIFJUTA: \{'C55', 'A43', 'A67', 'C43', 'C10', 'A18', 'C67'\}} \\
\\ 
\texttt{ARGUMENT: \{'Din il{-}proposta toħloq isforz totali ta’ 1433, iżda l{-}importanza totali tiegħi għal dawn l{-}oġġetti hija biss 3726, li hija baxxa ħafna. F’dan il{-}kuntest, l{-}oġġetti li joffruni l{-}aħjar ratio ta’ importanza għal isforz huma "B42" (783/390=2.01), "C73" (508/849=0.60), u "B36" (461/350=1.32). Nistgħu nikkonċentraw fuq sett li jinkludi oġġetti b’ratio għoli u jibqa’ fil{-}limitu.'\}} \\
\\ 
\texttt{PROPOSTA: \{'B42', 'C73', 'B36', 'C61', 'C42'\}} \\
            }
        }
    }
    & \\ \\

    \theutterance \stepcounter{utterance}  
    & \multicolumn{4}{p{0.6\linewidth}}{
        \cellcolor[rgb]{0.9,0.9,0.9}{
            \makecell[{{p{\linewidth}}}]{
                \texttt{\tiny{[P1$\rangle$GM]}}
                \texttt{L{-}isforz totali ta' din il{-}proposta hija: **390 (B42) + 849 (C73) + 350 (B36) + 630 (C61) + 801 (C42) = 3020** (li hija **< 3288**).} \\
\\ 
\texttt{L{-}importanza totali **tiegħi** hija: **783 (B42) + 508 (C73) + 461 (B36) + 668 (C61) + 808 (C42) = 3228**.} \\
\\ 
\texttt{**ARGUMENT:** \{'Din il{-}proposta toffri valur għoli għalija (3228), u għalkemm l{-}isforz qed jintuża b'mod effiċjenti (92\% tal{-}limitu), nista' nikkunsidra li naċċetta jekk tkun hemm tweġiba li tkompli toptimizza l{-}bilanċ bejn isforz u importanza. Jekk trid, nistgħu nikkonsidraw li nissostitwixxi oġġett b'isforz għoli b'oġġett ieħor li jkun ta' valur għoli għalija.'\}} \\
\\ 
\texttt{**Qed nistenna r{-}reazzjoni tiegħek.**} \\
            }
        }
    }
    & & \\ \\

    \theutterance \stepcounter{utterance}  
    & & & \multicolumn{2}{p{0.3\linewidth}}{
        \cellcolor[rgb]{0.9,0.9,0.9}{
            \makecell[{{p{\linewidth}}}]{
                \texttt{\tiny{[GM$|$GM]}}
                \texttt{abort} \\
            }
        }
    }
    & & \\ \\

    \theutterance \stepcounter{utterance}  
    & & & \multicolumn{2}{p{0.3\linewidth}}{
        \cellcolor[rgb]{0.9,0.9,0.9}{
            \makecell[{{p{\linewidth}}}]{
                \texttt{\tiny{[GM$|$GM]}}
                \texttt{end game} \\
            }
        }
    }
    & & \\ \\

\end{supertabular}
}

\end{document}
